\section{Cayley-Hamilton}
\label{section-cayley-hamilton}

\begin{lemma}
\label{lemma-charpoly}
Let $R$ be a ring. Let $A = (a_{ij})$ be an $n \times n$
matrix with coefficients in $R$. Let $P(x) \in R[x]$
be the characteristic polynomial of $A$ (defined
as $\det(x\text{id}_{n \times n} - A)$).
Then $P(A) = 0$ in $\text{Mat}(n \times n, R)$.
\end{lemma}

\begin{proof}
We reduce the question to the well-known Cayley-Hamilton
theorem from linear algebra in several steps:
\begin{enumerate}
\item If $\phi :S \rightarrow R$ is a ring morphism and $b_{ij}$
are inverse images of the $a_{ij}$ under this map, then it suffices
to show the statement for $S$ and $(b_{ij})$ since $\phi$ is a ring morphism.
\item If $\psi :R \hookrightarrow S$ is an injective ring morphism, it
clearly suffices to show the result for $S$ and the $a_{ij}$ considered as
elements of $S$. 
\item Thus we may first reduce to the case $R = \mathbf{Z}[X_{ij}]$,
$a_{ij} = X_{ij}$ of a polynomial ring and then further to
the case $R = \mathbf{Q}(X_{ij})$ where we may finally apply Cayley-Hamilton.
\end{enumerate}
\end{proof}

\begin{lemma}
\label{lemma-charpoly-module}
Let $R$ be a ring.
Let $M$ be a finite $R$-module.
Let $\varphi : M \to M$ be an endomorphism.
Then there exists a monic polynomial $P \in R[T]$ such that
$P(\varphi) = 0$ as an endomorphism of $M$.
\end{lemma}

\begin{proof}
Choose a surjective $R$-module map $R^{\oplus n} \to M$, given by
$(a_1, \ldots, a_n) \mapsto \sum a_ix_i$ for some generators $x_i \in M$.
Choose $(a_{i1}, \ldots, a_{in}) \in R^{\oplus n}$ such that
$\varphi(x_i) = \sum a_{ij} x_j$. In other words the diagram
$$
\xymatrix{
R^{\oplus n} \ar[d]_A \ar[r] & M \ar[d]^\varphi \\
R^{\oplus n} \ar[r] & M
}
$$
is commutative where $A = (a_{ij})$. By
Lemma \ref{lemma-charpoly}
there exists a monic polynomial $P$ such that $P(A) = 0$.
Then it follows that $P(\varphi) = 0$.
\end{proof}

\begin{lemma}
\label{lemma-charpoly-module-ideal}
Let $R$ be a ring. Let $I \subset R$ be an ideal.
Let $M$ be a finite $R$-module.
Let $\varphi : M \to M$ be an endomorphism such
that $\varphi(M) \subset IM$.
Then there exists a monic polynomial
$P = t^n + a_1 t^{n - 1} + \ldots + a_n \in R[T]$
such that $a_j \in I^j$ and $P(\varphi) = 0$ as an endomorphism of $M$.
\end{lemma}

\begin{proof}
Choose a surjective $R$-module map $R^{\oplus n} \to M$, given by
$(a_1, \ldots, a_n) \mapsto \sum a_ix_i$ for some generators $x_i \in M$.
Choose $(a_{i1}, \ldots, a_{in}) \in I^{\oplus n}$ such that
$\varphi(x_i) = \sum a_{ij} x_j$. In other words the diagram
$$
\xymatrix{
R^{\oplus n} \ar[d]_A \ar[r] & M \ar[d]^\varphi \\
I^{\oplus n} \ar[r] & M
}
$$
is commutative where $A = (a_{ij})$. By
Lemma \ref{lemma-charpoly}
the polynomial
$P(t) = \det(t\text{id}_{n \times n} - A)$
has all the desired properties.
\end{proof}

\noindent
As a fun example application we prove the following surprising lemma.

\begin{lemma}
\label{lemma-fun}
Let $R$ be a ring.
Let $M$ be a finite $R$-module.
Let $\varphi : M \to M$ be a surjective $R$-module map.
Then $\varphi$ is an isomorphism.
\end{lemma}

\begin{proof}[First proof]
Write $R' = R[x]$ and think of $M$ as a finite $R'$-module with
$x$ acting via $\varphi$. Set $I = (x) \subset R'$. By our assumption that
$\varphi$ is surjective we have $IM = M$. Hence we may apply
Lemma \ref{lemma-charpoly-module-ideal}
to $M$ as an $R'$-module, the ideal $I$ and the endomorphism $\text{id}_M$.
We conclude that
$(1 + a_1 + \ldots + a_n)\text{id}_M = 0$ with $a_j \in I$.
Write $a_j = b_j(x)x$ for some $b_j(x) \in R[x]$.
Translating back into $\varphi$ we see that
$\text{id}_M = -(\sum_{j = 1, \ldots, n} b_j(\varphi)) \varphi$, and hence
$\varphi$ is invertible.
\end{proof}

\begin{proof}[Second proof]
We perform induction on the number of generators of $M$ over $R$. If
$M$ is generated by one element, then $M \cong R/I$ for some ideal
$I \subset R$. In this case we may replace $R$ by $R/I$ so that $M = R$.
In this case $\varphi : R \to R$ is given by multiplication on $M$ by an
element $r \in R$. The surjectivity of $\varphi$ forces $r$ invertible,
since $\varphi$ must hit $1$, which implies that $\varphi$ is
invertible.

\medskip\noindent
Now assume that we have proven the lemma in the case of modules
generated by $n - 1$ elements, and are examining a module $M$ generated
by $n$ elements. Let $A$ mean the ring $R[t]$, and regard the module
$M$ as an $A$-module by letting $t$ act via $\varphi$; since $M$ is
finite over $R$, it is finite over $R[t]$ as well, and since we're
trying to prove $\varphi$ injective, a set-theoretic property, we might
as well prove the endomorphism $t : M \to M$ over $A$ injective. We have
reduced our problem to the case our endomorphism is multiplication by
an element of the ground ring. Let $M' \subset M$ denote the
sub-$A$-module generated by the first $n - 1$ of the generators of $M$,
and consider the diagram
$$
\xymatrix{
0 \ar[r] & M' \ar[r]\ar[d]^{\varphi\mid_{M'}} & M\ar[d]^\varphi \ar[r] &
M/M' \ar[d]^{\varphi \bmod M'} \ar[r] & 0 \\
0 \ar[r] & M' \ar[r]                                  & M \ar[r]
           & M/M' \ar[r]                                  & 0,
}
$$
where the restriction of $\varphi$ to $M'$ and the map induced by $\varphi$
on the quotient $M/M'$ are well-defined since $\varphi$ is multiplication
by an element in the base, and $M'$ and $M/M'$ are $A$-modules in
their own right. By the case $n = 1$ the map $M/M' \to M/M'$ is an
isomorphism. A diagram chase implies that $\varphi|_{M'}$ is surjective
hence by induction $\varphi|_{M'}$ is an isomorphism. This forces the
middle column to be an isomorphism by the snake lemma.
\end{proof}













